\documentclass[11pt]{article}
\usepackage[margin=1in]{geometry}
\usepackage{fontspec}
\usepackage{hyperref}
\usepackage{orcidlink}
\usepackage{longtable}
\usepackage{fvextra}
\fvset{breaklines=true,breakanywhere=true,fontsize=\small}
\title{Dossier de Registro: padic-grn-core\\\large Núcleo de análisis p-ádico para sistemas dinámicos discretos}
\author{
Jesús Rogelio Pérez-Buendía\,\orcidlink{0000-0002-7739-4779}
\and
Víctor Nopal-Coello\,\orcidlink{0000-0003-2608-3636}
\and
Yuriria Cortés-Poza\,\orcidlink{0000-0003-1856-9206}
}
\date{2026-04-25}

\begin{document}
\maketitle
\begin{abstract}
\texttt{padic-grn-core} es el núcleo matemático y computacional para calcular métricas p-ádicas (\texttt{mu\_E}, \texttt{mu\_A}, \texttt{mu\_I}) en redes de regulación genética y, con adaptaciones controladas, en sistemas dinámicos discretos finitos más generales. Este dossier reúne documentación funcional, código fuente, pruebas y evidencia de validación reproducible para fines de registro de software.
\end{abstract}
\tableofcontents
\newpage

\section{Introducción}
Este software centraliza el cálculo de perfiles p-ádicos de expansión, contracción e indiferencia sobre dinámicas discretas deterministas. Su diseño prioriza trazabilidad científica, reproducibilidad y reutilización por terceros: incluye interfaz de línea de comandos, adaptadores de entrada, pruebas automatizadas y ejemplos mínimos verificables.

El objetivo práctico de \texttt{padic-grn-core} es ofrecer una base confiable sobre la cual puedan construirse módulos de búsqueda (\texttt{padic-grn-search}) y de visualización (\texttt{padic-grn-viz}), manteniendo consistencia matemática y operativa en todo el ecosistema.

\section{Datos generales}
\begin{itemize}
  \item Nombre: \texttt{padic-grn-core}
  \item Versión: \texttt{0.1.0}
  \item Lenguaje: Python 3.10+
  \item Tipo: software de cálculo de métricas p-ádicas para sistemas discretos finitos
  \item DOI: \texttt{DOI\_SOFTWARE\_CORE}
\end{itemize}

\subsection{Correos electrónicos oficiales de autores}
\begin{itemize}
  \item Jesús Rogelio Pérez-Buendía: \texttt{rogelio.perez@cimat.mx}
  \item Víctor Nopal-Coello: \texttt{victor\_nopal@uaeh.edu.mx}
  \item Yuriria Cortés-Poza: \texttt{yuriria.cortes@iimas.unam.mx}
\end{itemize}

\subsection{Financiamiento}
\begin{itemize}
  \item Modelos matemáticos y computacionales no convencionales para el estudio y análisis de problemas relevantes en Biología, CF 2019/217367.
\end{itemize}

\section{Autoría oficial}
\VerbatimInput{../../../AUTORIA-ORCID-DOI-OFICIAL.md}

\section{Documentación funcional completa}
\subsection{README}
\VerbatimInput{../README.md}
\subsection{Manual de usuario}
\VerbatimInput{../MANUAL-USUARIO.md}
\subsection{Manual técnico}
\VerbatimInput{../MANUAL-TECNICO.md}
\subsection{Configuración de paquete}
\VerbatimInput{../pyproject.toml}
\subsection{Cita académica}
\VerbatimInput{../CITATION.cff}
\subsection{Licencia}
\VerbatimInput{../LICENSE}
\subsection{Configuración de pruebas}
\VerbatimInput{../pytest.ini}

\section{Código fuente completo}
\subsection{Inicialización de paquete}
\VerbatimInput{../src/padic_grn_core/__init__.py}
\subsection{CLI}
\VerbatimInput{../src/padic_grn_core/cli.py}
\subsection{Adaptadores}
\VerbatimInput{../src/padic_grn_core/adapters.py}
\subsection{Script de validación}
\VerbatimInput{../tools/validate_expected_values.py}

\section{Pruebas y ejemplos completos}
\subsection{Prueba smoke}
\VerbatimInput{../tests/test_smoke.py}
\subsection{Pruebas de adaptadores}
\VerbatimInput{../tests/test_adapters.py}
\subsection{Ejemplo de transiciones p=2}
\VerbatimInput{../examples/transitions_p2_n2_identity.csv}
\subsection{Ejemplo de transiciones p=3}
\VerbatimInput{../examples/transitions_p3_n2_identity.csv}
\subsection{Ejemplo de orderings}
\VerbatimInput{../examples/orderings_n2.csv}

\section{Evidencia textual de validación}
\subsection{Log de ejecución de pruebas}
\VerbatimInput{validacion_core_pytest.txt}
\subsection{Salida JSON toy p=2}
\VerbatimInput{../../../registro-dossier/06_evidencia_corridas/core_scores_p2.json}
\subsection{Salida JSON toy p=3}
\VerbatimInput{../../../registro-dossier/06_evidencia_corridas/core_scores_p3.json}
\subsection{Salida CSV batch}
\VerbatimInput{../../../registro-dossier/06_evidencia_corridas/core_batch_p2.csv}
\subsection{Salida A. thaliana}
\VerbatimInput{../../../registro-dossier/06_evidencia_corridas/core_scores_athaliana_pi_star.json}
\subsection{Salida C. elegans}
\VerbatimInput{../../../registro-dossier/06_evidencia_corridas/core_scores_celegans_pi_star.json}
\subsection{Salida Drosophila}
\VerbatimInput{../../../registro-dossier/06_evidencia_corridas/core_scores_drosophila_pi_star.json}

\end{document}
